\chapter{Internet of Things laboratorio}

\section{Pasarela}
\begin{itemize}
\item Como forma de tolerancia a fallos, usted decide poner tres pasarelas entre su red IoT e Internet.
  \begin{itemize}
  \item ¿Cómo evito que me llegen los mensajes por triplicado?
  \item ¿Necesito enviar tres mensajes distintos desde un sensor?
  \item ¿Cómo compruebo que todos los nodos están vivos?
  \end{itemize}
\end{itemize}

\section{MQTT}

\section{Topología}
Las redes de sensores inalámbricas son unas redes formadas por nodos de bajo coste dotados con una tecnología inalámbrica (por ejemplo 802.15.4e). Una de las configuraciones mas habituales de ese tipo de redes es la comunicación multisalto. Esta configuración es usada para extender el alcance de comunicación de cada uno de los nodos gracias a que cada uno de los nodos enruta la información del resto hacia la pasarela. La pasarela envía toda la información a través de un interfaz mas apropiado (por ejemplo, un modem LTE). El problema es que un nodo sólo puede enrutar tráfico estando en estado DESPIERTO ya que en estado DORMIDO no recibe ningún tipo de mensaje. De forma adicional, un nodo constantemente DESPIERTO consume mucha batería por lo que hay que mantener el mínimo número de nodos despiertos imprescindible para que la información pueda llegar a la pasarela. Se le pide que diseñe un algoritmo a ejecutar en cada uno de los nodos de la red, de forma que:

Al final de la ejecución del algoritmo, un nodo decida si está DESPIERTO o DORMIDO.
Un nodo DORMIDO se despierta regularmente y envía la lectura de su sensor (por ejemplo, temperatura) a un nodo qué, para que reciba su mensaje, tiene que estar  DESPIERTO.
Un nodo DESPIERTO tiene a todos sus vecinos dormidos excepto a uno que se encuentra mas cerca de la pasarela y le servirá para enviar la información a la pasarela.  Un vecino de un nodo es un nodo con el cual tiene enlace directo. 

En la siguiente figura vemos una red de sensores inalámbrica donde cada nodo viene representado por un círculo y cada enlace inalámbrico directo viene representado por una linea. Si cada uno de los nodos  ejecuta su algoritmo, los nodos tomarán una decisión de si se quedan despiertos o dormidos formando un camino hacia la pasarela. Por ejemplo, en la red de la figura una posible solución sería la que se representa en la figura de abajo con los nodos que han decidido quedarse despiertos en negro.
Como consejo para el desarrollo del algoritmo, comience con redes simples (dos o tres nodos) conectados a una pasarela y trate de dibujar a mano el intercambio de mensajes necesario para que esa red simple logre el objetivo planteado. Trate de extender ese primer algoritmo a redes mas complejas.


% Local Variables:
%   coding: utf-8
%   fill-column: 90
%   mode: flyspell
%   ispell-local-dictionary: "spanish"
%   mode: latex
%   TeX-master: "main"
% End:
